% Lösungen Einheit 1 von 5 – Matrizen als Rechenobjekte (zusammenbau v0.1)
\einheitenkopf{Einheit 1 von 5 · Matrizen als Rechenobjekte}

\erg{9}{a) rechnen \quad b) $A \cdot v = (4 | 11)$ \quad c) $A \cdot B$ ist bildbar mit Zeilenzahl $= 2$ und Spaltenzahl $= 4$; $B \cdot A$ ist nicht bildbar, weil B 4 Spalten, A aber nur 2 Zeilen hat \quad d) $M^2 = ((5 | 4), (4 | 5))$ \quad e) $A^{-1} = ((1 | -1), (-1 | 2))$ (aus $2a + c = 1$, $a + c = 0$, $2b + d = 0$, $b + d = 1$) \quad f) P vertauscht die beiden Komponenten: $P \cdot (a | b) = (b | a)$ \quad g) $a = 3$, $b = 4$ (aus $a + 2 = 5$ und $2 + 2b = 10$) \quad h) genau für $b = 0$ und $d = a$ (c beliebig): Ausmultiplizieren ohne Vertauschen gibt $X \cdot A = A \cdot X$, also $((b | 0), (d | 0)) = ((0 | 0), (a | b))$}
\erg{10}{a) Fehler: elementweise quadriert statt Zeile mal Spalte. Richtig: $M^2 = ((4 | 5), (0 | 9))$ \quad b) Beim Produkt wird jede Zeile des linken mit jeder Spalte des rechten Faktors kombiniert; vertauscht man die Faktoren, treffen andere Zeilen auf andere Spalten – die Ergebnisse stimmen im Allgemeinen nicht überein}
